\section{Cryptographic primitives and representation costs}
\label{sec:shuffle-primitives}

\subsection{Hashing a secret-shared 32-byte query key}
The query key is shared between the garbler and evaluator:
$x=x_G\oplus x_E\in\{0,1\}^{256}$.  Neither party is assumed to know $x$.
Consequently, computing its two table addresses locally is not an admissible
optimization.  A concrete baseline evaluates AES-128-CMAC inside the secure
computation, with a fixed, separately sampled secret PRF key $K$ shared between
the parties.  The database uses the same function when it is constructed.
For the fixed-length input $x=M_1\mathbin\Vert M_2$, CMAC is
\begin{align}
 L&=\operatorname{AES}_K(0^{128}), &
 K_1&=\operatorname{dbl}(L),\nonumber\\
 C_1&=\operatorname{AES}_K(M_1), &
 T&=\operatorname{AES}_K(C_1\oplus M_2\oplus K_1).
 \label{eq:shared-cmac}
\end{align}
Both message blocks are complete, so CMAC uses $K_1$, without a padding block.
Multiplication by $X$ in the binary field, denoted $\operatorname{dbl}$, consists
of a shift and XORs with selected bits of the constant reduction polynomial;
it requires no AND gates.  These choices follow AES and the CMAC specification
\cite{aesfips,cmac}.

For an auditable gate count, use the established 32-AND AES S-box circuit listed
by NIST, rather than assuming a newer optimized implementation.  The ten AES
rounds contain $160$ S-boxes, giving $160\cdot32=5120$ ANDs for the encryption
datapath.  The AES-128 key schedule contains another $40$ S-boxes, or $1280$
ANDs.  NIST lists $6400$ ANDs for AES-128 with its key schedule and $5120$ for
the fixed-key datapath \cite{nistcircuits}.  A secret, already expanded key has
the same nonlinear datapath count: AddRoundKey remains XOR.  Reusing the
expanded key and $K_1$ across the $B$ queries of a single batch therefore gives
\begin{equation}
 A_{\rm hash}(B)=
 \underbrace{1280+5120}_{6400\text{ setup}}
 +\underbrace{2\cdot5120B}_{10240B\text{ query work}}.
 \label{eq:shared-hash-cost}
\end{equation}
At 32 bytes per half-gate AND, this is $200$ KiB of setup plus $320$ KiB per
query.  The primary model retains the 11 expanded round keys and $K_1$ as
1536 ordinary encoded state bits, using the same linked-garbling interface
as other persistent state.  Setup is therefore performed once per initialized
state and amortized, rather than repeated in every batch.  This requires
compatible labels and distinct gate/hash tweaks across the composed
computation; it does not permit reuse of one-time garbled tables.  A reset to
an incompatible encoding must pay for secure key-state translation or new
input and setup.  Equation~\eqref{eq:shared-hash-cost} counts neither initial
database hashing nor interface traffic.

One 128-bit PRF result suffices for two $h$-bit table addresses when $2h\le128$:
take two disjoint $h$-bit slices of $T$.  For a pseudorandom $T$, the slices are
jointly pseudorandom; a second CMAC is unnecessary for this address-generation
interface.  Equal keys deliberately receive equal address pairs.  The full
query key must still be compared with candidate keys to handle address
collisions.  This construction assumes fixed-length, domain-separated use of
the PRF; changing its application domain must preserve that separation.

Query-share input and result-share output belong in a separate interface ledger.
Reconstructing $x_G\oplus x_E$ inside GC costs XORs but not ANDs.  The model
charges 4096 bytes for G's share labels and 12,288 bytes for E's input OT,
giving 16,384 bytes per query: 12,288 forward and 4096 reverse, before framing.
This is a convention, not an unavoidable lower bound: existing garbled inputs
and correlated-OT encodings can change the charge.  The value and hit flag
remain shared; taking the two parties' label color bits gives local XOR
shares as explained in Section~\ref{sec:interactive-pool-protocol}.
There is no plaintext answer decoding to E.  Client delivery is a separate
boundary and is not included in the inter-server totals.

\subsection{Packed storage and conditional movement are separate costs}
The proposed packed representation carries a 32-byte value in two 128-bit
objects.  We retain this as a payload-storage and transport optimization.
It does not, by itself, specify the protocol for moving those objects under
a secret control bit.  In ordinary Boolean GC, a 128-bit label represents
one Boolean wire; storing a payload in two opaque carriers is a different
representation with its own input, output, and movement interfaces.

For this reason the primary total uses a fully specified Boolean movement
cost: swapping a 256-bit payload under one secret bit costs 256 ANDs, with
metadata added at its actual Boolean width.  A conditional packed-backend
column may instead use a measured or proved parameter $c_2$, the communication
for swapping the two packed payload carriers under a secret control.  Storage
size alone does not establish $c_2$, and the conditional column must also
include conversions and relabeling.  When a shuffle-and-reveal step has made
the physical selection public, both parties can select their own shares
locally; that public routing genuinely avoids secret-control payload gates.
This distinction makes the Boolean estimate auditable while keeping the
packed optimization visible as a concrete implementation target.

\subsection{One-hot garbling: what is obtained and what is charged}
Let $P=2^k$ be the padded size of a small array.  Heath's binary-to-one-hot
primitive creates $H(y)$ using $(k-1)\lambda$ bits of garbled material, or
$16(k-1)$ bytes at $\lambda=128$, when the evaluator knows $y$.
The complete framework is not restricted to public true inputs: it masks
private inputs and cancels the masks within secure modules
\cite{heath2021,heath2024}.

In particular, Section~5.1 of \cite{heath2024} encodes a private short integer
$x\in\mathbb Z_{2^k}$ using a mask $\alpha$ and a one-hot encoding of
$y=x+\alpha\bmod 2^k$.  G knows $\alpha$ and E learns $y$, while neither learns
$x$.  For binary input, the stated construction adds $\alpha$ with a
$(k-1)$-AND ripple-carry adder and then invokes binary-to-one-hot.  With the
paper's two-ciphertext Boolean AND implementation, the conversion material is
\begin{equation}
 C_{\rm bin\to masked\ hot}
 =32(k-1)+16(k-1)=48(k-1)\ \text{bytes}.
 \label{eq:masked-hot}
\end{equation}
This excludes input-label transfer, the small reveal/interface bits, and any
additional conversion required by the consuming operation.  Binary and
arithmetic one-hot outputs must not silently be identified: an arithmetic
consumer may require the corresponding upcast or word conversion.  The
expanded vector still has $P$ entries, so local work and working storage grow
linearly in $P$ for the binary construction, with further width factors for
arithmetic words or payload outer products.

Equation~\eqref{eq:masked-hot} is a cost for a \emph{masked-position}
representation, not a canonical $H(x)$ at the same interface.  A lookup into
a mutable secret array must also align that representation with the array
and account for any masking correction.  The free truth-table evaluation
$T(f)^\mathsf{T}H(y)$ uses a public constant table; it is not automatically a
lookup into independently encoded secret array entries.  Similarly, the cheap
one-hot outer product distributes one scalar or vector to one selected slot,
not a different unknown payload at every slot.  These observations do not
rule out a useful module, especially with interaction and shuffled arrays;
they identify the missing interface that must be specified.  The primary
totals therefore claim no generic one-hot discount for dynamic pool accesses.

\subsection{One-AND swaps and a fused flat-array exchange}
For a control bit $s$ and two $w$-bit words $a,b$, define
\begin{equation}
 d=s\mathbin{\wedge}(a\oplus b),\qquad
 a'=a\oplus d,\qquad b'=b\oplus d.
 \label{eq:one-and-swap}
\end{equation}
The AND is bitwise with the same control on all bits.  Thus a swap costs
$w$ ANDs and $3w$ XORs, rather than two independently implemented muxes.
For $P$ flat entries and a valid unique target index, start a carry word with
the new value and conditionally swap it with each array entry at its one-hot
selection bit.  The final carry is the old value and the array contains the
new one.  A straightforward shared-prefix binary decoder uses $P-2$ ANDs
for $P\ge2$, giving the specified circuit cost
\begin{equation}
 A_{\rm flat\ exchange}(P,w)=Pw+P-2,
 \qquad X_{\rm flat\ exchange}(P,w)=3Pw+P-2.
 \label{eq:flat-fused-exchange}
\end{equation}
These are counts for this explicit decoder, not a claim of minimum circuit
size.  A one-entry array needs no selector circuit.  An additional secret
enable, out-of-range handling, or metadata predicates must be charged if
present.  The fused exchange assumes the replacement word is already
available; an arbitrary function of the returned old word needs its own
evaluation and suitable writeback.  For batched repeated indices, a
single unordered aggregate of selectors is insufficient: preserve sequential
semantics or securely group duplicates and define their read/write order.
